% ============================================================
%  N.1.2  Workloads
% ============================================================
\subsection{Workloads}
\label{sec:workloads}

Two classes of CNN workloads were selected for the experiments:
\emph{activation-dominant} and \emph{weight-dominant}
(Table~\ref{tab:workloads}).
The rationale for including both classes is that layer fusion
primarily eliminates expensive DRAM reads and writes of intermediate
activations, which would normally be transferred to and from external
memory in single-layer or layer-by-layer processing
(see Section~\ref{sec:partial-full-fusion}).
Activation-dominant workloads, characterised by large spatial dimensions
and small channel counts, should therefore benefit the most from fusion,
because intermediate activations constitute the dominant fraction of
off-chip traffic.
It is equally important, however, to evaluate fusion on weight-dominant
workloads, where large channel counts and shrinking spatial dimensions
cause weight-related DRAM traffic to outweigh activation traffic,
making the gains from eliminating intermediate transfers less
straightforward.

Two of the four workloads, VGG16 and ResNet18, contain stride-2
downsampling layers.
When fusing across stride boundaries, the input tile size of a
subsequent layer must account for the cumulative stride of all
preceding layers (e.g.\ a $7 \times 7$ output tile at the end of
ResNet18 requires a $112 \times 112$ input tile at the first layer),
which increases the on-chip memory pressure for the fused workload.

% ------------------------------------------------------------------
%  Workload summary table
% ------------------------------------------------------------------
\begin{table}[htbp]
  \centering
  \caption{Summary of the four CNN workloads used in the experiments.}
  \label{tab:workloads}
  \small
  \begin{tabular}{l c c c c c}
    \toprule
    \textbf{Workload}
      & \textbf{Type}
      & \textbf{Fused layers}
      & \textbf{Input size}
      & \textbf{Max channels}
      & \textbf{Stride} \\
    \midrule
    FSRCNN~\cite{fsrcnn_paper}
      & Activation-dom.
      & 8
      & $960 \times 540$
      & 56
      & None \\
    MC-CNN~\cite{mccnn_paper}
      & Activation-dom.
      & 4
      & $1242 \times 376$
      & 32
      & None \\
    VGG16~\cite{vgg_paper}
      & Weight-dom.
      & 13
      & $224 \times 224$
      & 512
      & $2$ (4 layers) \\
    ResNet18~\cite{resnet_paper}
      & Weight-dom.
      & 17
      & $224 \times 224$
      & 512
      & $2$ (4 layers) \\
    \bottomrule
  \end{tabular}
\end{table}

\paragraph{FSRCNN}~\cite{fsrcnn_paper} is a super-resolution network
composed of 8~convolutional layers with mixed filter sizes
($5 \times 5$, $1 \times 1$, $3 \times 3$).
The channel count remains small throughout (maximum~56), while the
spatial resolution stays constant at $960 \times 540$ with no stride
layers, making intermediate activation volumes very large relative to
weight volumes.

\paragraph{MC-CNN}~\cite{mccnn_paper} is a stereo-matching network with
4~homogeneous convolutional layers, all using $3 \times 3$ filters and
a uniform channel count of~32.
The input resolution of $1242 \times 376$ is the largest among the four
workloads, making it strongly activation-dominant.

\paragraph{VGG16}~\cite{vgg_paper} comprises 13~convolutional layers,
all with $3 \times 3$ filters, organised in 5~blocks.
Channels grow from 64 to 512 while the spatial resolution shrinks from
$224 \times 224$ to $14 \times 14$ through four stride-2 transitions
(modelling the max-pool between blocks), yielding a cumulative stride
of~16.
Three fully connected layers (FC6--FC8) are excluded from fusion, as
they exhibit different data access patterns.

\paragraph{ResNet18}~\cite{resnet_paper} contains 17~main convolutional
layers (excluding 3~shortcut projection $1 \times 1$ convolutions and
the final FC layer, which are computed separately).
The first layer uses a $7 \times 7$ filter with stride~2; three
additional stride-2 layers occur at stages~3, 4, and~5, for a
cumulative stride of~16.
Channels increase from 64 to 512 while the spatial resolution
decreases from $56 \times 56$ (after the first layer) to $7 \times 7$.
